% META: {"id": "1NSI-PM-COURS-04", "chapitre": "1NSI-PROJET-METHODES", "type_objet": "cours", "capacites": ["1NSI-PROJET-METHODES-C4"], "mode_creation": "ex_nihilo", "sources_inspiration": [], "status": "approved", "capacites_codes": ["C4"]}
\section{Documenter son code et préparer une présentation orale}

\definition[D4]{
  \textbf{Documenter} un programme, c'est expliquer, en langage naturel, ce qu'il fait :
  le rôle de chaque fonction (via une \textbf{docstring}), les choix de conception
  importants, et la manière de l'utiliser.
}

\begin{codereference}{Une fonction bien documentée}
\begin{python}
def calcul_moyenne(valeurs):
    """Calcule la moyenne arithmetique d'une liste de valeurs numeriques.

    Precondition : valeurs est une liste non vide.
    Postcondition : le resultat est compris entre min(valeurs) et max(valeurs).
    """
    return sum(valeurs) / len(valeurs)

def a_une_docstring(fonction):
    """Teste si `fonction` possede une docstring non vide."""
    return fonction.__doc__ is not None and len(fonction.__doc__.strip()) > 0

print(a_une_docstring(calcul_moyenne))  # True
\end{python}
\end{codereference}

% BEGIN-VERIFY
% def calcul_moyenne(valeurs):
%     """Calcule la moyenne arithmetique d'une liste de valeurs numeriques."""
%     return sum(valeurs) / len(valeurs)
% def a_une_docstring(fonction):
%     return fonction.__doc__ is not None and len(fonction.__doc__.strip()) > 0
% assert a_une_docstring(calcul_moyenne) is True
% def sans_doc(x):
%     return x
% assert a_une_docstring(sans_doc) is False
% END-VERIFY

\propriete[Ce qu'une bonne documentation contient]{
  Une docstring de fonction utile précise, dans l'ordre : \textbf{ce que fait} la fonction
  (une phrase claire), ses \textbf{préconditions} (ce que les arguments doivent vérifier),
  et éventuellement ses \textbf{postconditions} (ce que garantit le résultat) — exactement
  la spécification vue au chapitre sur les langages et la programmation.
}

\propriete[Préparer une présentation orale de projet]{
  Le programme officiel souligne que cette spécialité contribue au développement des
  \textbf{compétences orales}, en particulier pour préparer l'épreuve orale de terminale
  (Grand Oral) — un travail qui commence dès la Première. Une présentation de projet
  efficace explique : le \textbf{problème} traité, les \textbf{choix} effectués (et
  pourquoi), les \textbf{difficultés} rencontrées et comment elles ont été surmontées, et
  une \textbf{démonstration} concrète du résultat.
}

\erreurFrequente{Préparer une présentation orale qui ne fait que \textbf{décrire le code
ligne par ligne}. Un jury (ou un public) s'intéresse davantage au \textbf{raisonnement} :
pourquoi ce problème est intéressant, pourquoi telle approche a été choisie plutôt qu'une
autre, ce qui a été appris en le résolvant — pas à une lecture littérale du code source.}

\alamachine{Choisis une fonction que tu as écrite dans un chapitre précédent de ce manuel
(par exemple le tri par insertion ou la recherche dichotomique). Écris-lui une docstring
complète (rôle, précondition, postcondition), puis vérifie avec \lstinline{a_une_docstring}
qu'elle est bien documentée. Prépare ensuite, à l'oral, une explication de $2$ minutes de
cette fonction pour un camarade qui ne l'a jamais vue.}
